\documentclass[11pt]{article}
\usepackage[a4paper,margin=25mm]{geometry}
\usepackage{amsmath,amssymb,amsthm}
\usepackage[hidelinks]{hyperref}
\newtheorem{theorem}{Theorem}
\title{Algebraic Multiplicity Does Not Determine the Perturbation Exponent:\\A Disproof of Conjecture 00000005131}
\author{GitHub: gaochengzhecpu}
\date{}
\begin{document}
\maketitle
\begin{abstract}
The zero matrix of order two has an eigenvalue of algebraic multiplicity
two, but its worst perturbational eigenvalue displacement is exactly
$\varepsilon$, and its pseudospectral area is exactly
$\pi\varepsilon^2$. These are not the proposed sharp powers
$\varepsilon^{1/2}$ and $\varepsilon^{2/2}$. Lean verifies the actual
complex matrices, all admissible perturbations, the characteristic
polynomial and multiplicity, the exact pseudospectrum, its Lebesgue
area, and the negations of the corresponding universal power laws.
\end{abstract}

\paragraph{Reading of the source.}
Both source languages explicitly introduce $m$ as \emph{algebraic
multiplicity}, assert displacement of $\varepsilon^{1/m}$ type, and
assert pseudospectral area of $\varepsilon^{2/m}$ type. We interpret a
stated scaling exponent as a sharp power: comparability up to positive
constants, or a positive finite leading asymptotic constant. We refute
even the lower bound necessary for either meaning. A mere big-$O$
upper bound is a weaker statement and is not refuted here.

The same source says the exponent is decided by the maximal Jordan
block size. That is the correct distinction, and it conflicts with
using algebraic multiplicity as the exponent denominator. Our result
does \emph{not} refute the version in which $m$ is consistently replaced
by maximal Jordan block size. In the example that size is one. The
independent, explicit area clause also fails with $m=2$, so this is
not an argument based only on one chosen perturbation path.

\paragraph{Precise conventions.}
For a complex matrix $E$ use the induced infinity norm
\[
 \|E\|_\infty=\max_i\sum_j |E_{ij}|.
\]
Its closed perturbational pseudospectrum is
\[
 \Sigma_\varepsilon(A)=
 \{z\in\mathbb C:\text{some }E\text{ satisfies }
 \|E\|_\infty\le\varepsilon\text{ and }(A+E)v=zv
 \text{ for some }v\ne0\}.
\]
Area means ordinary planar Lebesgue measure. The example has only one
spectral cluster, avoiding any distinction between local and total
pseudospectral area. Open perturbation balls give the same area and
the same supremal displacement.

\begin{theorem}
Let $A=0\in M_2(\mathbb C)$ and $\varepsilon\ge0$. Then
\[
 \chi_A(X)=X^2,\qquad
 \Sigma_\varepsilon(A)=\{z:|z|\le\varepsilon\},\qquad
 \operatorname{area}(\Sigma_\varepsilon(A))=\pi\varepsilon^2.
\]
In particular the eigenvalue zero has algebraic multiplicity two,
and the maximum possible eigenvalue displacement is exactly
$\varepsilon$.
\end{theorem}
\begin{proof}
The characteristic polynomial is $\det(XI_2)=X^2$, so the root
multiplicity of zero is two. Let $Ev=zv$ with $v\ne0$, and choose a
coordinate $i$ with $|v_i|=\max_j|v_j|>0$. Every perturbation in the
stated ball satisfies
\[
 |z|\,|v_i|
 =\left|\sum_jE_{ij}v_j\right|
 \le\sum_j|E_{ij}|\,|v_j|
 \le\left(\sum_j|E_{ij}|\right)|v_i|
 \le\varepsilon|v_i|.
\]
Thus $|z|\le\varepsilon$. Conversely, if $|z|\le\varepsilon$,
take $E=zI_2$. Its row sums are $|z|$, and it has eigenvalue $z$
with nonzero eigenvector $(1,1)$. This proves the exact disk
description, including attainment of displacement $\varepsilon$ at
$z=\varepsilon$. Its Lebesgue area is $\pi\varepsilon^2$.
\end{proof}

\begin{theorem}
Neither the sharp displacement exponent $1/m$ nor the sharp
pseudospectral area exponent $2/m$ is universally determined by
algebraic multiplicity $m$.
\end{theorem}
\begin{proof}
Use the preceding matrix, with $m=2$. If the claimed displacement
exponent were sharp, some $c,\delta>0$ would imply
$c\sqrt\varepsilon\le\varepsilon$ whenever
$0<\varepsilon<\delta$. Dividing by $\sqrt\varepsilon$ would give
$c\le\sqrt\varepsilon$ at arbitrarily small positive
$\varepsilon$, a contradiction. Similarly, the area law would require
$c\varepsilon\le\pi\varepsilon^2$ on a positive interval, which would
give $c\le\pi\varepsilon$ there, again impossible.

The Lean statements use $\varepsilon=t^m$ and $t\downarrow0$ to
express both necessary lower bounds without fractional powers. The
response lower bound is $ct$ and the area lower bound is $ct^2$.
For the example the exact values are $t^2$ and $\pi t^4$,
respectively. All quantifiers over positive constants and all
sufficiently small parameters are retained in the proofs.
\end{proof}

\paragraph{Norm and Jordan-block scope.}
The same counterexample applies to the customary Euclidean operator
norm: the eigenvector estimate $|z|\|v\|\le\|E\|\|v\|$ and
$\|zI_2\|=|z|$ give the identical disk. Lean specifically verifies
the induced infinity norm version above; this Euclidean norm
observation is a paper proof. Both are standard fixed norms, and
finite-dimensional equivalence of norms cannot change these sharp
exponents. The zero matrix is diagonal and has maximal Jordan block
size one, exactly consistent with the corrected exponent. No claim
is made about a different corrected Jordan-block formulation or
the source's condition-number and lattice-configuration clauses.
Falsity of the universal area clause with algebraic $m$ already
refutes the stated conjunction.

\paragraph{Formal verification and scope.}
The Lean development uses Mathlib matrices over $\mathbb C$, actual
matrix-vector multiplication, Mathlib characteristic polynomials
and polynomial root multiplicity, and Mathlib planar volume. The
predicate \texttt{Budget} is the row-sum characterization of the
induced infinity norm ball; \texttt{Eigen} is the ordinary nonzero
eigenvector equation. The exact disk theorem quantifies over
\emph{all} matrices in that ball. The principal results are
\begin{center}
\texttt{zero\_pseudospectrum}, \texttt{sharp\_response},
\texttt{zero\_multiplicity},\\
\texttt{zero\_area}, \texttt{not\_multiplicity\_area\_law},\\
\texttt{not\_multiplicity\_response\_law}.
\end{center}
The area law is universally quantified over single-cluster matrices
with $\chi_A=(X-a)^n$, so algebraic multiplicity is $n$ and no
unrelated eigenvalue can contribute extra area. It states only the
positive lower bound required by the source's area scaling claim;
negating that necessary consequence suffices. The response law
similarly includes actual characteristic-polynomial multiplicity.
Only the standard axioms \texttt{propext}, \texttt{Classical.choice}
and \texttt{Quot.sound} occur. Development and self-review were
performed by a delegated AI agent; the coordinating agent reviews
the result separately. No external independent review is claimed.

\paragraph{Source and reproduction.}
The preserved English and Chinese source is \texttt{SOURCE.md}, from
the repository's conjecture 00000005131. The provenance file records
the upstream commit and source hash. The Lean project pins
Lean 4.19.0 and Mathlib commit
\begin{center}\texttt{c44e0c8ee63ca166450922a373c7409c5d26b00b}.\end{center}
From the \texttt{lean} directory run
\begin{center}\texttt{lake build}\\
\texttt{lake env lean -DwarningAsError=true Main.lean}.\end{center}
The source ends with axiom reports for the principal theorems.
\end{document}
